\documentclass[10pt,a4paper]{amsart}
\usepackage[margin=19mm]{geometry}
\usepackage{amsmath,amssymb,mathtools}
\usepackage[scr=rsfs,cal=euler]{mathalfa}
\usepackage{xfrac}
\usepackage[unicode,hidelinks,bookmarks=false]{hyperref}
\newcommand{\ie}{i.e.\ }
\newcommand{\cf}{cf.\ }
\newcommand{\resp}{respectively\ }
\newcommand{\Subsection}{\subsection}
\providecommand{\nobreakdash}{\nobreak\hbox{-}}
\setlength{\emergencystretch}{2em}
\multlinegap0pt
\usepackage{bm}
\usepackage{bbm}
\usepackage{dsfont}
\usepackage{xspace}
\usepackage{cancel}
\usepackage{mathtools}
\usepackage{amsthm}
\makeatletter
\@ifundefined{proposition}{\newtheorem{proposition}{Proposition}}{}
\@ifundefined{remark}{\newtheorem{remark}{Remark}}{}
\makeatother
\title[From Dirac to Dunkl Operators through Symmetry Reduction]{From Dirac to Dunkl Operators through Symmetry Reduction}
\author{Cristina Sard\'on}
\date{Source: arXiv:2510.08283v1 (CC BY 4.0)}
\begin{document}
\maketitle

\noindent\emph{Source and licence.} From Dirac to Dunkl Operators through Symmetry Reduction. Version arXiv:2510.08283v1 (CC BY 4.0), distributed under the Creative Commons Attribution 4.0 International licence, as recorded in the arXiv licence metadata for this exact version and shown on the abstract page https://arxiv.org/abs/2510.08283v1. Full rights notice: licenses/arxiv-2510.08283-CC-BY-4.0.txt.\par\smallskip

\noindent\textit{Editorial scope.} Counted unit: the Proposition whose source label is \texttt{prop1} in the article (source0.tex lines 324--337, source label \texttt{prop1}), delivered together with the article's own complete original proof of it (source0.tex lines 339--427, one \texttt{proof} environment of 89 lines). No mathematics is written, repaired or re-derived: statement and proof are the article's own text line for line. Required context retained uncounted: none. Every definition, display and label of those lines is retained unchanged. Omitted from this package and not redistributed: 1--323, 338--338, 428--1501. Nothing inside a retained excerpt is omitted or altered: every retained body is delivered exactly as the article's lines are written, so no omission receipt of any kind accompanies this package. Citations: the retained excerpts carry no citation command at all, so no bibliography entry of the article is transcribed and no reference is redistributed. Reference adaptation: none was needed and none was made; every reference inside the delivered unit resolves inside it, so no unresolved reference or citation is produced. Printed slips kept literal: no printed word, symbol, hypothesis or number of the counted statement or of its proof was altered. Layout and class: the article's own class is \texttt{article} with packages including inputenc, fontenc, lmodern, amsmath, amssymb, amsthm, mathtools, bm, bbm, mathrsfs, geometry, enumitem, hyperref, authblk; the wrapper is the lane's pinned \texttt{amsart} editorial wrapper at 10pt on A4, which restates the theorem-like environments the retained text needs and adds the article's own macro definitions that the retained text needs (none), with extra wrapper packages (bm, bbm, dsfont, xspace, cancel, mathtools). The delivered numbering is the unit's own. Assets: no retained span contains a figure environment, a graphics-inclusion command, a PDF-page inclusion, a table or a TikZ picture; no image file is copied into this package, the package declares no asset dependency and no image file is redistributed with it. Licence: version arXiv:2510.08283v1 of this article is distributed under the Creative Commons Attribution 4.0 International licence, as recorded in the arXiv licence metadata for this exact version and shown on the abstract page https://arxiv.org/abs/2510.08283v1. A full rights notice is delivered at licenses/arxiv-2510.08283-CC-BY-4.0.txt. Characters: every retained excerpt is pure ASCII, so the delivered engine, XeLaTeX with its default Latin Modern text font, reports no missing character.
\begin{proposition}\label{prop1}
For $\xi \in V$, the adjoint of the directional derivative
$\partial_\xi$ with respect to the weighted inner product
$\langle \cdot , \cdot \rangle_k$ is given by
\[
  \partial_\xi^{*} \;=\; -\partial_\xi \;-\; \langle \nabla \log \delta_k(x), \xi \rangle ,
\]
ignoring boundary terms. In particular,
\[
  \nabla \log \delta_k(x) = \sum_{\alpha \in R^+} 
  \frac{2k(\alpha)}{\langle \alpha, x\rangle}\,\alpha .
\]

\end{proposition}

\begin{proof}
Considering the weighted inner product
\[
  \langle f,g \rangle_k
  \;=\;
  \int_C f(x)\,\overline{g(x)}\,\delta_k(x)\,dx .
\]
\noindent
Integrating by parts in this weighted measure, the derivative does not only move across, but it is picked on the weight \(\delta_k(x)\), i.e., 

\[
  \int f(x)\,\overline{\partial_\xi g(x)}\,\delta_k(x)\,dx  = -\int (\partial_\xi f)(x)\,\overline{g(x)}\,\delta_k(x)\,dx
    \;-\; \int f(x)\,\overline{g(x)}\,(\partial_\xi \delta_k(x))\,dx .
\]



Now divide the second term by \(\delta_k(x)\) on the right-hand side to write it as
\[
  - \int f(x)\,\overline{g(x)}\,
     \frac{\partial_\xi \delta_k(x)}{\delta_k(x)}\,\delta_k(x)\,dx .
\]

But
\begin{equation}\label{division}
  \frac{\partial_\xi \delta_k(x)}{\delta_k(x)}
  \;=\; \langle \nabla \log \delta_k(x), \xi \rangle.
\end{equation}

The previous expression comes from the fact that the directional derivative of a smooth function $f$ in the direction $\xi$ is defined by $\partial_\xi f(x) := \langle \nabla f(x), \xi \rangle$. Applying this to $f(x) = \log \delta_k(x)$ and using the chain rule  we retrieve \eqref{division}.



So altogether:
\[
  \langle f, \partial_\xi g \rangle_k
  \;=\;
  \langle \big(-\partial_\xi - \langle \nabla \log \delta_k(x), \xi \rangle\big) f , g \rangle_k .
\]

which is the definition of the adjoint \(\partial_\xi^*\):
\begin{equation}\label{defadjoint}
  \partial_\xi^* \;=\; -\partial_\xi \;-\; \langle \nabla \log \delta_k(x), \xi \rangle .
\end{equation}
See that one retrieves $\partial_{\xi}^{*} \;=\; -\partial_{\xi}$, the plain Lebesgue inner product, by canceling out the term on the right-hand side involving the special measure $\delta_k(x).$


\begin{remark}
The above calculation is to be understood in the sense that we ignore potential boundary contributions arising from
integration by parts.  For functions with compact support inside $C$
such terms vanish, and the formula for $\partial_\xi^{*}$ is \eqref{defadjoint}.
For more general domains of definition, one has to specify suitable
boundary conditions to obtain a rigorous operator-theoretic adjoint.
\end{remark}

We define the drift term:
\begin{equation}\label{driftdef}
  \nabla \log \delta_{k}(x)
  = \sum_{\alpha \in R^{+}} \frac{2k(\alpha)}{\langle \alpha,x\rangle}\, \alpha .
\end{equation}
\noindent
which appears from the following computation on the weight
\[
  \delta_k(x) \;=\; \prod_{\alpha \in R^+} |\langle \alpha , x \rangle|^{2k(\alpha)} .
\]

Taking logarithm on both sides of the previous expression, 
\[
  \log \delta_k(x)
  \;=\; \sum_{\alpha \in R^+} 2k(\alpha)\, \log |\langle \alpha , x \rangle| 
\]
\noindent
and further differentiation on both sides (gradient) using that \(\nabla \langle \alpha , x \rangle = \alpha\) provides:
\[
  \nabla \log \delta_k(x)
  \;=\; \sum_{\alpha \in R^+} \frac{2k(\alpha)}{\langle \alpha , x \rangle}\, \alpha .
\]


Again, with the definition of  the directional derivative
\(\partial_\xi := \langle \nabla , \xi \rangle\),  the correction term in the adjoint formula is
\[
  \langle \nabla \log \delta_k(x), \xi \rangle
  \;=\; \sum_{\alpha \in R^+} 
         \frac{2k(\alpha)\langle \alpha , \xi \rangle}{\langle \alpha , x \rangle}.
\]


\end{proof}

\end{document}
